%!TEX root = ../dokumentation.tex
\thispagestyle{plain}

\chapter*{Abkürzungsverzeichnis}
\markboth{Abkürzungsverzeichnis}{Abkürzungsverzeichnis}
%nur verwendete Akronyme werden letztlich im Dokument angezeigt
\begin{acronym}[MOSFET]
    %\setlength{\itemsep}{-\parsep}
    %Manuell nach Alphabet sortieren
    \acro{ATL}{Automatic-Tools-List (dt. Automatisierte-Tool-Liste)}
    \acro{CI}{Continuous Integration}
    \acro{C++}{C Plus Plus  (Programmiersprache)}
    \acro{CS}{C Sharp (Programmiersprache)}
    \acro{DLL}{Dynamic Link Library}
    \acro{LFS}{Large File Storage}
    \acro{LIN}{LinInducitonsSimulator2 (Intern geschriebenes Programm für die Simulation von Induktionserwärmung bestimmter Heizplatten)}
    \acro{NWA}{Nutzwertanalyse}
    \acro{SVN}{Subversion}
    \acro{URL}{Uniform Resource Locator}
    \acro{VCS}{Version Control System}
    
    %
    \acroplural{DLL}[DLLs]{Dynamic Link Libraries}
    \acroplural{VCS}[VCSs]{Version Control Systems (dt. Versionsverwaltungssysteme)}
    %Manuell nach Alphabet sortieren
    %VS Code -> Zeilen markieren STRG + Shift + P, nach Sort Lines Ascending suchen
\end{acronym}